% chapters/appendix_a.tex
\section*{ДОДАТОК А}
\addcontentsline{toc}{section}{ДОДАТОК А. Фрагменти вихідного коду симулятора}
{\centering\bfseries Додаток А\par}
{\centering Фрагменти вихідного коду програмної реалізації\par}
\vspace{1em}

% -------------------------------------------------------------------
\subsection*{А.1 Визначення типів даних (src/core/types.ts)}
% -------------------------------------------------------------------

\begin{lstlisting}[style=jsstyle,
  caption={Визначення типів даних симулятора},
  label={lst:types}]
// src/core/types.ts

export type NodeId = string;
export type PortId = string;

export type Complex = { re: number; im: number };

export type Token = {
  id: string;
  value: number | Complex;
  t: number;       // logical arrival time
  originT: number;
};

export type Edge = {
  id: string;
  from: { node: NodeId; port: PortId };
  to:   { node: NodeId; port: PortId };
  delay?: number;      // channel delay (simulation ms)
  label?: string;      // arc label text
  labelValue?: Complex; // numeric value (tooltip)
};

export type NodeSpec = {
  id: NodeId;
  kind: NodeKind;
  inPorts: PortId[];
  outPorts: PortId[];
  latency: number;  // computation time (simulation ms)
  state?: Record<string, unknown>;
  params?: {
    twiddle?: { N: number; k: number };
  };
};

export type Graph = {
  nodes: NodeSpec[];
  edges: Edge[];
};

export type FireResult = {
  outputs: Record<PortId, Token>;
};

export type NodeKind =
  'source' | 'sink' | 'add' | 'mul' | 'butterfly' | 'dft4' | 'twiddle';
\end{lstlisting}

% -------------------------------------------------------------------
\subsection*{А.2 Функції комплексної арифметики (src/utils/math.ts)}
% -------------------------------------------------------------------

\begin{lstlisting}[style=jsstyle,
  caption={Функції комплексної арифметики та алгоритм 4-точкового ДПФ},
  label={lst:math}]
// src/utils/math.ts
import { Complex } from '../core/types';

export const cplxAdd = (a: Complex, b: Complex): Complex => ({
  re: a.re + b.re,
  im: a.im + b.im,
});

export const cplxSub = (a: Complex, b: Complex): Complex => ({
  re: a.re - b.re,
  im: a.im - b.im,
});

// Multiply by j: (re + im*j) * j = -im + re*j
export const cplxMulJ = (a: Complex): Complex => ({
  re: -a.im,
  im: a.re,
});

// Optimised 4-point DFT without non-trivial multiplications
export function computeDFT4(
  x0: Complex, x1: Complex, x2: Complex, x3: Complex
): Complex[] {
  // E0 = x0 + x2
  const E0 = cplxAdd(x0, x2);
  // E1 = x0 - x2
  const E1 = cplxSub(x0, x2);
  // O0 = x1 + x3
  const O0 = cplxAdd(x1, x3);
  // O1 = x1 - x3
  const O1 = cplxSub(x1, x3);

  // X0 = E0 + O0
  const X0 = cplxAdd(E0, O0);
  // X2 = E0 - O0
  const X2 = cplxSub(E0, O0);

  // X1 = E1 - j*O1
  const jO1 = cplxMulJ(O1);
  const X1  = cplxSub(E1, jO1);
  // X3 = E1 + j*O1
  const X3  = cplxAdd(E1, jO1);

  return [X0, X1, X2, X3];
}
\end{lstlisting}

% -------------------------------------------------------------------
\subsection*{А.3 Генерація графу потоку даних (src/graph/generateDFT4x4.ts)}
% -------------------------------------------------------------------

\begin{lstlisting}[style=jsstyle,
  caption={Генерація DFG обчислювача 16-точкового ШПФ},
  label={lst:generate}]
// src/graph/generateDFT4x4.ts

import type { Graph, NodeSpec, Edge } from '@/core/types';

export function generateDFT4x4(latency = 400): Graph {
  const nodes: NodeSpec[] = [];
  const edges: Edge[] = [];

  const N = 16;
  const radix = 4;

  // 1. CREATE SOURCE NODES
  for (let i = 0; i < N; i++) {
    nodes.push({
      id: `src${i}`,
      kind: 'source',
      inPorts: [],
      outPorts: ['out'],
      latency: 0,
    });
  }

  // 2. CREATE ROW DFT4 NODES
  for (let n1 = 0; n1 < radix; n1++) {
    nodes.push({
      id: `row-${n1}`,
      kind: 'dft4',
      inPorts: ['in0', 'in1', 'in2', 'in3'],
      outPorts: ['out0', 'out1', 'out2', 'out3'],
      latency,
    });

    // Connect Source -> Row DFT
    // Input array index: i = n1 + 4*n2
    for (let n2 = 0; n2 < radix; n2++) {
      const srcIndex = n1 + 4 * n2;
      edges.push({
        id: `edge-src${srcIndex}-row${n1}in${n2}`,
        from: { node: `src${srcIndex}`, port: 'out' },
        to:   { node: `row-${n1}`,      port: `in${n2}` },
        delay: 600,
      });
    }
  }

  // 3. CREATE TWIDDLE FACTOR NODES  W_16^(n1 * k2)
  for (let n1 = 0; n1 < radix; n1++) {
    for (let k2 = 0; k2 < radix; k2++) {
      const twiddlePower = n1 * k2;
      const twiddleId    = `tw-${n1}-${k2}`;

      nodes.push({
        id: twiddleId,
        kind: 'twiddle',
        inPorts:  ['in'],
        outPorts: ['out'],
        latency: latency / 2, // multiplication is faster than DFT4
        params: {
          twiddle: { N: 16, k: twiddlePower },
        },
      });

      // Connect Row DFT -> Twiddle
      edges.push({
        id: `edge-row${n1}out${k2}-tw${n1}_${k2}`,
        from:  { node: `row-${n1}`, port: `out${k2}` },
        to:    { node: twiddleId,   port: 'in' },
        label: twiddlePower > 0 ? `W16_${twiddlePower}` : undefined,
        delay: 600,
      });
    }
  }

  // 4. CREATE COLUMN DFT4 NODES
  for (let k2 = 0; k2 < radix; k2++) {
    nodes.push({
      id: `col-${k2}`,
      kind: 'dft4',
      inPorts:  ['in0', 'in1', 'in2', 'in3'],
      outPorts: ['out0', 'out1', 'out2', 'out3'],
      latency,
    });

    // Connect Twiddle -> Column DFT (matrix transpose)
    // Column node k2 collects outputs from all row nodes n1
    for (let n1 = 0; n1 < radix; n1++) {
      const twiddleId = `tw-${n1}-${k2}`;
      edges.push({
        id:   `edge-tw${n1}_${k2}-col${k2}in${n1}`,
        from: { node: twiddleId,  port: 'out' },
        to:   { node: `col-${k2}`, port: `in${n1}` },
        delay: 600,
      });
    }
  }

  // 5. CREATE SINK NODES
  for (let i = 0; i < N; i++) {
    nodes.push({
      id: `snk${i}`,
      kind: 'sink',
      inPorts:  ['in'],
      outPorts: [],
      latency:  0,
    });
  }

  // Connect Column DFT -> Sink
  // Output index: k = k2 + 4*k1
  for (let k2 = 0; k2 < radix; k2++) {
    for (let k1 = 0; k1 < radix; k1++) {
      const sinkIndex = k2 + 4 * k1;
      edges.push({
        id:   `edge-col${k2}out${k1}-sink${sinkIndex}`,
        from: { node: `col-${k2}`,     port: `out${k1}` },
        to:   { node: `snk${sinkIndex}`, port: 'in' },
        delay: 600,
      });
    }
  }

  return { nodes, edges };
}
\end{lstlisting}

% -------------------------------------------------------------------
\subsection*{А.4 Короткий опис модуля ядра симуляції (src/core/scheduler.ts)}
% -------------------------------------------------------------------

Модуль \texttt{src/core/scheduler.ts} реалізує клас \texttt{DataflowRuntime} ---
ядро виконання dataflow-моделі.
Конструктор приймає об'єкт \texttt{Graph} (вузли та дуги) і функцію
\texttt{getTime(): number}, що повертає поточний логічний час симуляції.
При ініціалізації для кожної дуги графа створюється окрема FIFO-черга
типу \texttt{Map<edgeId, Token[]>}, а для кожного вузла --- запис у
\texttt{Map<nodeId, number>} (час готовності \texttt{readyAt}).

Основний метод \texttt{tick(dt: number, maxFires: number)} виконує чотири
кроки на кожному кроці симуляції:
\begin{enumerate}
  \item \textbf{Доставка токенів.} Для кожної дуги перевіряється умова
    $t_{\text{now}} - t_{\text{token}} \geq \mathtt{edge.delay}$; токени, що
    задовольняють умову, переміщуються з черги дуги до вхідного буфера
    \texttt{inbox} відповідного вузла.
  \item \textbf{Перевірка готовності вузла.} Вузол вважається готовим до
    виконання, якщо всі його вхідні порти заповнені токенами і поточний
    логічний час не менший за \texttt{readyAt[nodeId]}.
  \item \textbf{Активація вузла.} Для готового вузла викликається метод
    \texttt{execute()}, що обчислює вихідні токени залежно від типу вузла
    (\texttt{NodeKind}): \texttt{dft4} викликає \texttt{computeDFT4()},
    \texttt{twiddle} --- комплексне множення на $W_N^k$, \texttt{sink} ---
    зберігає результат у масиві виходів.
    Значення \texttt{readyAt[nodeId]} оновлюється на
    $t_{\text{now}} + \mathtt{node.latency}$.
  \item \textbf{Постановка вихідних токенів у черги.} Метод
    \texttt{emitFrom()} поміщає кожен вихідний токен у FIFO-черги дуг,
    що виходять з активованого вузла, з міткою часу $t_{\text{emit}}$.
\end{enumerate}

Клас підтримує механізм подій: \texttt{onToken} (токен доставлено на вхід
вузла), \texttt{onFire} (вузол активовано), \texttt{onOutput} (результат
з'явився на виході \texttt{sink}).
Ці події використовуються компонентом \texttt{GraphView} для анімації
руху токенів і оновлення теплової карти активності вузлів.

Повний код модуля наведено у репозиторії проекту.
